% Syntra LaTeX Template
% For Arabic support, use XeLaTeX with:
%   \usepackage{fontspec}
%   \usepackage{polyglossia}
%   \setdefaultlanguage{english}
%   \setotherlanguage{arabic}
%   \newfontfamily\arabicfont[Script=Arabic]{Amiri}

\documentclass[12pt,a4paper]{article}

% Font and encoding
\usepackage[utf8]{inputenc}
\usepackage[T1]{fontenc}

% Math packages
\usepackage{amsmath,amssymb}

% Table support
\usepackage{longtable,booktabs}

% Graphics and links
\usepackage{graphicx}
\usepackage{hyperref}
\usepackage{geometry}

% Page geometry
\geometry{
    left=2.5cm,
    right=2.5cm,
    top=3cm,
    bottom=3cm
}

% Fix pandoc tightlist issue
\providecommand{\tightlist}{\setlength{\itemsep}{0pt}\setlength{\parskip}{0pt}}

% Hyperref settings
\hypersetup{
    colorlinks=true,
    linkcolor=blue,
    filecolor=magenta,
    urlcolor=cyan,
    pdftitle={Syntra Research Documentation},
    pdfauthor={Syntra-Env},
    pdfpagemode=FullScreen
}

\begin{document}

\section{Syntra: Technical
Implementation}\label{syntra-technical-implementation}

The Prime Framework \& Holonomic Dynamics · UOR v3.5.0 Specification ·
March 2026

\begin{quote}
Syntra aligns with \textbf{UOR v3.5.0}, implementing a \textbf{Prime
Framework} for universal number representation. This ensures
\textbf{Observer Invariance} and \textbf{Geometric Consistency} by
encoding information through complete prime factorization.
\end{quote}

\subsection{1. Universal Coordinate
Projection}\label{universal-coordinate-projection}

Content is projected onto a triadic UOR manifold using prime-factorized
Universal Numbers. Similarity is measured via \textbf{Geometric
Congruence} across prime spectra.

\subsubsection{Projection Strategy}\label{projection-strategy}

Every textual unit is mapped to an address where coordinates are stored
as \texttt{\{prime:\ exponent\}} factorizations. This eliminates the
"floating point drift" typical of standard vector embeddings.

\begin{verbatim}
# Conceptual UOR Projection
from uor import triality

c1 = triality.project_coordinate("Hello world")
c2 = triality.project_coordinate("Goodbye world")

# Similarity is 'resonance' between prime factors
resonance = triality.compute_resonance(c1, c2)
print(f"Universal Resonance Fidelity: {resonance:.4f}")
\end{verbatim}

\subsection{2. Holonomic Fields \& Semantic
Tension}\label{holonomic-fields-semantic-tension}

The \textbf{Holonomic Field} utilizes prime factorization to compute
observer-invariant congruence. \textbf{Semantic Tension}
(\$\textbackslash mathcal\{T\}\$) measures local curvature within the
prime-factorized semantic lattice.

\subsubsection{Manifold Proximity}\label{manifold-proximity}

Proximity is calculated using prime-aware distance metrics, ensuring
that related concepts share a high degree of algebraic congruence.

\begin{verbatim}
# Semantic Proximity on the Manifold
from uor import triality, field

dataset = ["AI", "Machine Learning", "Neural Network", "Recipe"]
manifold_points = [{'label': s, 'coordinate': triality.project_coordinate(s)} for s in dataset]

query = triality.project_coordinate("Deep Learning training")
neighbors = field.manifold_proximity(query, manifold_points, k=3)
\end{verbatim}

\subsection{3. Gauge Stability \&
Topology}\label{gauge-stability-topology}

The topology module assesses \textbf{Gauge Stability} by measuring
resonance fidelity across edges. This ensures mathematical consistency
across the manifold as the agent moves through the text (Parallel
Transport).

\begin{verbatim}
# Gauge Stability Check
from uor import topology

topology = topology.build_topology(nodes, edges)
stability = topology.assess_gauge_stability(topology)

print(f"Gauge Stability: {stability['resonance_fidelity']:.4f}")
\end{verbatim}

\subsection{4. Architectural Guarantees}\label{architectural-guarantees}

\begin{itemize}
\tightlist
\item
  \textbf{Observer Invariance:} The prime-factorized identity of a
  "glyph" is the same regardless of the context it is viewed in.
\item
  \textbf{Geometric Consistency:} All operations on the manifold follow
  the \$\textbackslash mathfrak\{su\}(2)\$ gauge field laws, providing
  hard mathematical limits on "meaning drift."
\end{itemize}

\end{document}
